\section{Coquand--Huet 1988: indexed constructions and the full beta calculus}
\label{sec:p07-ch88}

Coquand and Huet's journal presentation uses one written universe symbol,
indexed binding, and separate validity and typing judgments. Its logical
``contexts'' are themselves product terms; they must not be identified
with a modern list-of-declarations judgment or with an additional universe
constant. The presentation below follows the source's concrete binding
semantics in an editorially corrected indexed reconstruction. It separates
the full calculus from the corrected restricted variant and supplies a
normalization-preserving interpretation into the independently proved
two-sort CC. The original article and the separate 1994 normalization proof
retain their own attribution
\cite[\S\S1--3, pp.~95--105]{CoquandHuet1988}
\cite[\S\S2--3]{Geuvers1994Report9450}.

\subsection{One-star syntax and binding}

\begin{hbdefinition}{def:p07-ch88-syntax}{Objects and logical contexts at a fixed depth}
For \(n\geq0\), define mutually the scoped source classes
\(\mathcal A_i^n\) of objects, \(\mathcal A_o^n\) of logical contexts,
and \(\mathcal A^n=\mathcal A_i^n\cup\mathcal A_o^n\).
A logical context is \(\ast\), or \([x:A]\Delta\) with
\(A\in\mathcal A^n\), \(\Delta\in\mathcal A_o^{n+1}\).
An object is an index \(k\) with \(1\leq k\leq n\), a product
\([x:A]M\), or an abstraction \((\lambda x:A)M\), with
\(A\in\mathcal A^n\), \(M\in\mathcal A_i^{n+1}\); it may also
be \((MN)\) with \(M,N\in\mathcal A_i^n\). The two classes
are disjoint. Products and lambdas bind index one in the body, not in
their annotation. The written \(x\) is a dummy concrete binder name;
the abstract structure uses positive de Bruijn indices. Under \(d\)
internal binders, indices at most \(d\) refer to those binders and an
index \(k>d\) denotes external reference \(k-d\).

There is no written \(\Box\), and \(\ast\) has no source type.
The word object here includes proposition schemas and types, not just
the inhabitants called objects in the later two-sort classification.
No primitive eta, delta, inductive or recursor computation is included
in this base syntax \cite[pp.~95--98]{CoquandHuet1988}.
\end{hbdefinition}

\begin{hbdefinition}{def:p07-ch88-environments}{The two roles of source contexts}
A source environment \(\Gamma\) is a closed logical-context term,
viewed as an ordered sequence of declarations. Define its source length
\(|\ast|=0\), \(|[x:A]\Delta|=1+|\Delta|\).
Here \(|\cdot|\) is the source's local sequence-length notation.
For \(\Gamma\) of length \(n\) and \(\Delta\in\mathcal A_o^n\),
\(\Gamma;\Delta\) replaces the final \(\ast\) of \(\Gamma\)
by \(\Delta\), with its already scoped external references. Put
\(\Gamma[x:A]=\Gamma;[x:A]\ast\).
Write
\[
 \Gamma=[x_n:A_n]\cdots[x_1:A_1]\ast,
 \qquad A_k\in\mathcal A^{n-k}.
\]
Thus index one names the newest declaration, whereas the leftmost
declaration is oldest. A logical context \(\Delta\) in a type position
is still a product term, not an instruction to append its declarations.
Validity will be \(\Gamma\vdash\Delta\); object typing will be
\(\Gamma\vdash M:P\). Their empty environment is the source term
\(\ast\), so \(\vdash E\) abbreviates \(\ast\vdash E\)
\cite[pp.~97--99]{CoquandHuet1988}.
\end{hbdefinition}

\begin{hbdefinition}{def:p07-ch88-index-operations}{Relocation, lookup and binder-removal substitution}
For cutoff \(i\geq0\), relocation \(\xi_i\) leaves \(\ast\)
unchanged and acts on an index by
\[
 \xi_i(k)=\begin{cases}k,&k\leq i,\\k+1,&k>i.\end{cases}
\]
It acts componentwise on application; on either binder it applies
\(\xi_i\) to the annotation and \(\xi_{i+1}\) to the body.
Consequently the already bound indices are fixed. Write
\(M^+=\xi_0(M)\), \(\operatorname{add}(0,M)=M\), and
\(\operatorname{add}(j+1,M)=\operatorname{add}(j,M)^+\).
For \(1\leq k\leq|\Gamma|\), the source lookup is
\(\Gamma/k=\operatorname{add}(k,A_k)\).

For an object argument \(N\in\mathcal A_i^n\), define
\(\operatorname{sub}_i(N,Q)\) recursively by fixing \(\ast\),
distributing over application, and applying cutoff \(i\) to a binder
annotation and cutoff \(i+1\) to its body. On indices it is
\[
 \operatorname{sub}_i(N,k)=
 \begin{cases}
 k,&k\leq i,\\
 \operatorname{add}(i,N),&k=i+1,\\
 k-1,&k>i+1.
 \end{cases}
\]
The source notation \([N/x]Q\) means
\(\operatorname{sub}_0(N,Q)\): replace the newest external variable
and remove its binder. Only \(N\) is inserted; the substituted image
is not recursively substituted into itself. The object-argument restriction
is necessary for the source's application grammar: arbitrary logical-context
terms are not raw application arguments. The relocation convention realizes
the source's free-reference-only binding semantics
\cite[pp.~98--99]{CoquandHuet1988}.
\end{hbdefinition}

\begin{hblemma}{thm:p07-ch88-index-operations}{proved=annex:p7.ch88.index-operations}{Scopedness of the actual index operations}
Relocation preserves the object/logical-context class and sends depth
\(n\) to depth \(n+1\). Lookup \(\Gamma/k\) is scoped at depth
\(|\Gamma|\). If \(N\in\mathcal A_i^n\) and
\(Q\in\mathcal A_b^{n+1}\), for \(b=i\) or \(b=o\), then
\([N/x]Q\in\mathcal A_b^n\). These are total structural operations
on their stated finite inputs; no normalization premise is needed.
\end{hblemma}

\begin{hbexample}{def:p07-ch88-index-example}{Lookup and capture avoidance}
In \(\Gamma=[a:\ast][t:1]\ast\), the newest index denotes \(t\),
and \(\Gamma/1=2\) denotes its older type \(a\), while
\(\Gamma/2=\ast\). In the extension \(\Gamma[x:2]\), take
\(Q=(\lambda z:3)2\): its annotation refers to \(a\), and its
body refers to \(x\), not the new \(z\). Substitution of the old
argument \(1\), denoting \(t\), gives
\([1/x]Q=(\lambda z:2)2\). The two occurrences of index two have
different scopes: the annotation refers to \(a\), and the body to
\(t\). No inserted reference has been captured by \(z\).
\end{hbexample}

\begin{hbdefinition}{def:p07-ch88-beta}{Full compatible annotated beta reduction}
The relation \(\longrightarrow_\beta\) is generated by the scoped
root contraction on Sheet~\ref{sheet:p07-ch88-beta}, closed under
both application children and both the annotation and body of a
product or lambda. The body's external index depth increases by one.
It is a relation on raw scoped objects and logical-context terms;
it retains annotations and has no restricted logical typing guard.
Write \(\longrightarrow_\beta^*\) for its reflexive-transitive
closure and \(\longrightarrow_\beta^+\) for its positive transitive
closure. This directed relation is not either variant's typed
conversion judgment \cite[p.~102]{CoquandHuet1988}.
\end{hbdefinition}

\subsection{Validity, typing and the two conversion regimes}

\begin{hbdefinition}{def:p07-ch88-judgments}{Common constructions and source-bound equality variants}
For an environment of length \(n\), a validity judgment
\(\Gamma\vdash\Delta\) has \(\Delta\in\mathcal A_o^n\);
a typing judgment \(\Gamma\vdash M:P\) has
\(M\in\mathcal A_i^n\), \(P\in\mathcal A^n\).
In particular \(\Gamma\vdash\ast\) says that the environment
is valid; it does not type \(\ast\) as its own inhabitant.
The common validity, product, projection, abstraction and application
rules are Sheet~\ref{sheet:p07-ch88-core}.

Typed conversion is written \(\Gamma\vdash M\cong N\).
The corrected restricted variant uses the source's logical guards:
reflexivity on valid logical contexts and on typed objects whose type
is a logical context; logical lambda and application compatibility;
and beta at a body typed by a logical context. The full variant uses
ordinary reflexivity on every typed object, arbitrary typed lambda
compatibility, simultaneous function/argument compatibility, and beta
at an arbitrary typed body. Both retain symmetry, transitivity, product
compatibility and typing conversion. Their separate sheets retain these
guards; neither typed equality is defined merely by removing a proof
annotation or by identifying every type with its inhabitants
\cite[\S3.1, pp.~101--102]{CoquandHuet1988}.
\end{hbdefinition}

\begin{hblemma}{thm:p07-ch88-restricted-inclusion}{proved=annex:p7.ch88.restricted-inclusion}{The corrected restricted variant embeds in the full variant}
Every validity, typing or typed-conversion judgment derivable in the
corrected restricted variant is derivable in the full variant, with
the same indexed syntax and environment. This is an inclusion of
derivations, not equality of the two conversion relations or an
identification of their type-checking procedures.
\end{hblemma}

\begin{hbdefinition}{def:p07-ch88-logical-reduction}{The source directed logical relation}
Write \(\Gamma\vdash M\mathrel{\triangleright}N\) for the corrected
source logical-reduction relation on
Sheet~\ref{sheet:p07-ch88-logical-reduction}. It has logical root-beta
and compatible annotation/body/application clauses, with the source's
logical typing guards. It also has an explicit transitivity rule.
A derived edge may therefore represent several contractions; it is
not a reflexive reduction relation and not exactly one ordinary beta step.
The source's logical rank is
\(\delta(M)=0\) for objects, \(\delta(\ast)=1\), and
\(\delta([x:M]\Delta)=\max(\delta(M)+1,\delta(\Delta))\)
for logical contexts. This rank is invariant under object substitution
and typed equality, as proved in \cref{thm:p07-ch88-check-rank}.
The domain-rank multiset built from it is nevertheless not decreasing
under innermost logical contraction
(\cref{thm:p07-ch88-check-domain-rank}). Logical termination is supplied
by the positive-path and annotated-normalization arguments below,
not by that multiset. The rank pertains to this logical regime,
not to an unrestricted source checking algorithm
\cite[pp.~104--105]{CoquandHuet1988}.
\end{hbdefinition}

\calculussheet{sheet:p07-ch88-core}{Coquand--Huet 1988 common constructions}{One-star scoped indexed objects and logical-context terms; separate validity and typing; corrected free-reference relocation/lookup/substitution; no type for star and no eta/delta or inductive extension.}
\hbrule{rule:p07-ch88-empty}{Empty validity}{}{\vdash\ast}{The empty environment is the logical-context term $\ast$.}
\hbrule{rule:p07-ch88-kind-extension}{Logical-context extension}{\Gamma\vdash\Delta}{\Gamma[x:\Delta]\vdash\ast}{$\Delta\in\mathcal A_o^{|\Gamma|}$.}
\hbrule{rule:p07-ch88-type-extension}{Object-type extension}{\Gamma\vdash M:\ast}{\Gamma[x:M]\vdash\ast}{$M\in\mathcal A_i^{|\Gamma|}$.}
\hbrule{rule:p07-ch88-kind-product}{Logical product}{\Gamma[x:M]\vdash\Delta}{\Gamma\vdash[x:M]\Delta}{The annotation is outside the new binder; the body is in its extended scope.}
\hbrule{rule:p07-ch88-type-product}{Object product}{\Gamma[x:M_1]\vdash M_2:\ast}{\Gamma\vdash[x:M_1]M_2:\ast}{The same source scoping convention.}
\hbrule{rule:p07-ch88-variable}{Projection}{\Gamma\vdash\ast}{\Gamma\vdash l:\Gamma/l}{$1\le l\le|\Gamma|$; lookup uses the corrected free-reference relocation.}
\hbrule{rule:p07-ch88-abstraction}{Abstraction}{\Gamma[x:M_1]\vdash M_2:P}{\Gamma\vdash(\lambda x:M_1)M_2:[x:M_1]P}{All fields have their declared indexed scope.}
\hbrule{rule:p07-ch88-application}{Application}{\Gamma\vdash M:[x:P]Q\qquad\Gamma\vdash N:P}{\Gamma\vdash(MN):[N/x]Q}{The capture-avoiding operation removes index one; the argument is an object.}

\calculussheet{sheet:p07-ch88-beta}{Full scoped annotated beta}{Raw indexed syntax with compatible closure through every application and binder annotation/body; distinct from the transitive typed logical relation.}
\hbrule{rule:p07-ch88-beta-root}{Root beta}{}{((\lambda x:A)M)N\longrightarrow_\beta[N/x]M}{$A\in\mathcal A^n$, $M\in\mathcal A_i^{n+1}$, $N\in\mathcal A_i^n$; close through both application children and both binder fields at their stated scopes.}

\calculussheet{sheet:p07-ch88-restricted}{Corrected source restricted conversion}{All common constructions; typed logical conversion only; ordinary reflexivity in the corrected reconstruction; logical guards retained; not the literal misprinted indexed surface.}
\hbrule{rule:p07-ch88-r-kind-refl}{Context reflexivity}{\Gamma\vdash\Delta}{\Gamma\vdash\Delta\cong\Delta}{\none}
\hbrule{rule:p07-ch88-r-refl}{Logical reflexivity}{\Gamma\vdash M:\Delta}{\Gamma\vdash M\cong M}{$\Delta$ is a logical-context term.}
\hbrule{rule:p07-ch88-r-sym}{Symmetry}{\Gamma\vdash M\cong N}{\Gamma\vdash N\cong M}{\none}
\hbrule{rule:p07-ch88-r-trans}{Transitivity}{\Gamma\vdash M\cong N\qquad\Gamma\vdash N\cong P}{\Gamma\vdash M\cong P}{\none}
\hbrule{rule:p07-ch88-r-product}{Product congruence}{\Gamma\vdash P_1\cong P_2\qquad\Gamma[x:P_1]\vdash M_1\cong M_2}{\Gamma\vdash[x:P_1]M_1\cong[x:P_2]M_2}{The body premise is in the first annotation's environment.}
\hbrule{rule:p07-ch88-r-abstraction}{Logical abstraction congruence}{\begin{gathered}\Gamma\vdash P_1\cong P_2\qquad\Gamma[x:P_1]\vdash M_1\cong M_2\\\Gamma[x:P_1]\vdash M_1:\Delta\end{gathered}}{\Gamma\vdash(\lambda x:P_1)M_1\cong(\lambda x:P_2)M_2}{$\Delta$ is a logical context.}
\hbrule{rule:p07-ch88-r-app-left}{Logical function congruence}{\Gamma\vdash(MN):\Delta\qquad\Gamma\vdash M\cong M_1}{\Gamma\vdash(MN)\cong(M_1N)}{$\Delta$ is a logical context.}
\hbrule{rule:p07-ch88-r-app-right}{Logical argument congruence}{\Gamma\vdash(MN):\Delta\qquad\Gamma\vdash N\cong N_1}{\Gamma\vdash(MN)\cong(MN_1)}{$\Delta$ is a logical context.}
\hbrule{rule:p07-ch88-r-beta}{Logical beta}{\Gamma[x:P]\vdash M:\Delta\qquad\Gamma\vdash N:P}{\Gamma\vdash((\lambda x:P)MN)\cong[N/x]M}{$\Delta$ is a logical context; the redex notation associates application to the whole lambda.}
\hbrule{rule:p07-ch88-r-conversion}{Typing conversion}{\Gamma\vdash M:P\qquad\Gamma\vdash P\cong Q}{\Gamma\vdash M:Q}{This is the source's typed equality premise, not an unqualified untyped equality test.}

\calculussheet{sheet:p07-ch88-full}{Coquand--Huet full conversion with corrected binding}{All common constructions; context reflexivity, symmetry/transitivity and product congruence as below; full replacements for the logical starred guards; no proof erasure or extra computation law.}
\hbrule{rule:p07-ch88-f-kind-refl}{Context reflexivity}{\Gamma\vdash\Delta}{\Gamma\vdash\Delta\cong\Delta}{\none}
\hbrule{rule:p07-ch88-f-refl}{Typed reflexivity}{\Gamma\vdash M:N}{\Gamma\vdash M\cong M}{Any source object typing; this is the printed full replacement.}
\hbrule{rule:p07-ch88-f-sym}{Symmetry}{\Gamma\vdash M\cong N}{\Gamma\vdash N\cong M}{\none}
\hbrule{rule:p07-ch88-f-trans}{Transitivity}{\Gamma\vdash M\cong N\qquad\Gamma\vdash N\cong P}{\Gamma\vdash M\cong P}{\none}
\hbrule{rule:p07-ch88-f-product}{Product congruence}{\Gamma\vdash P_1\cong P_2\qquad\Gamma[x:P_1]\vdash M_1\cong M_2}{\Gamma\vdash[x:P_1]M_1\cong[x:P_2]M_2}{The old binder environment is retained in the premise.}
\hbrule{rule:p07-ch88-f-abstraction}{Full abstraction congruence}{\begin{gathered}\Gamma\vdash P_1\cong P_2\qquad\Gamma[x:P_1]\vdash M_1\cong M_2\\\Gamma[x:P_1]\vdash M_1:N_1\end{gathered}}{\Gamma\vdash(\lambda x:P_1)M_1\cong(\lambda x:P_2)M_2}{The body type $N_1$ is arbitrary; no logical-context guard is imposed.}
\hbrule{rule:p07-ch88-f-application}{Full application congruence}{\begin{gathered}\Gamma\vdash(MN):P\qquad\Gamma\vdash M\cong M_1\\\Gamma\vdash N\cong N_1\end{gathered}}{\Gamma\vdash(MN)\cong(M_1N_1)}{Both children are replaced; the result type $P$ is arbitrary.}
\hbrule{rule:p07-ch88-f-beta}{Full typed beta}{\Gamma[x:A]\vdash M:P\qquad\Gamma\vdash N:A}{\Gamma\vdash((\lambda x:A)MN)\cong[N/x]M}{The body type $P$ is arbitrary.}
\hbrule{rule:p07-ch88-f-conversion}{Typing conversion}{\Gamma\vdash M:P\qquad\Gamma\vdash P\cong Q}{\Gamma\vdash M:Q}{No source-star self-type axiom is added.}

\calculussheet{sheet:p07-ch88-logical-reduction}{Corrected directed logical reduction}{Source page104 logical relation; compatible annotation/body steps with logical typing guards; explicit transitivity; no reflexive edge and no eta/delta rule.}
\hbrule{rule:p07-ch88-l-trans}{Path transitivity}{\Gamma\vdash M\mathrel{\triangleright}N\qquad\Gamma\vdash N\mathrel{\triangleright}P}{\Gamma\vdash M\mathrel{\triangleright}P}{A derived edge may comprise multiple contractions.}
\hbrule{rule:p07-ch88-l-product-body}{Product body}{\Gamma[x:P]\vdash M\mathrel{\triangleright}N}{\Gamma\vdash[x:P]M\mathrel{\triangleright}[x:P]N}{\none}
\hbrule{rule:p07-ch88-l-lambda-body}{Logical lambda body}{\Gamma[x:P]\vdash M\mathrel{\triangleright}N\qquad\Gamma[x:P]\vdash M:\Delta}{\Gamma\vdash(\lambda x:P)M\mathrel{\triangleright}(\lambda x:P)N}{$\Delta$ is a logical context.}
\hbrule{rule:p07-ch88-l-type-domain}{Object product annotation}{\Gamma[x:P]\vdash M:\ast\qquad\Gamma\vdash P\mathrel{\triangleright}Q}{\Gamma\vdash[x:P]M\mathrel{\triangleright}[x:Q]M}{The old body judgment supplies the source guard.}
\hbrule{rule:p07-ch88-l-lambda-domain}{Logical lambda annotation}{\Gamma[x:P]\vdash M:\Delta\qquad\Gamma\vdash P\mathrel{\triangleright}Q}{\Gamma\vdash(\lambda x:P)M\mathrel{\triangleright}(\lambda x:Q)M}{$\Delta$ is a logical context.}
\hbrule{rule:p07-ch88-l-kind-domain}{Logical product annotation}{\Gamma[x:P]\vdash\Delta\qquad\Gamma\vdash P\mathrel{\triangleright}Q}{\Gamma\vdash[x:P]\Delta\mathrel{\triangleright}[x:Q]\Delta}{\none}
\hbrule{rule:p07-ch88-l-app-left}{Logical function step}{\Gamma\vdash(MN):\Delta\qquad\Gamma\vdash M\mathrel{\triangleright}M_1}{\Gamma\vdash(MN)\mathrel{\triangleright}(M_1N)}{$\Delta$ is a logical context.}
\hbrule{rule:p07-ch88-l-app-right}{Logical argument step}{\Gamma\vdash(MN):\Delta\qquad\Gamma\vdash N\mathrel{\triangleright}N_1}{\Gamma\vdash(MN)\mathrel{\triangleright}(MN_1)}{$\Delta$ is a logical context.}
\hbrule{rule:p07-ch88-l-beta}{Logical root beta}{\Gamma[x:P]\vdash M:\Delta\qquad\Gamma\vdash N:P}{\Gamma\vdash((\lambda x:P)MN)\mathrel{\triangleright}[N/x]M}{$\Delta$ is a logical context.}


\subsection{The actual bridge to the proved two-sort CC}

\begin{hbdefinition}{def:p07-ch88-named-encoding}{A role-sensitive named encoding}
Fix an ordered external vector \(\vec z=(z_1,\ldots,z_n)\), oldest
to newest. Put \(T_{\vec z}(\ast)=\ast\) and
\(T_{\vec z}(k)=z_{n+1-k}\); preserve application, and map
\([x:A]B\) to \(\Pi z:T_{\vec z}(A).T_{\vec z,z}(B)\),
and \((\lambda x:A)M\) to the corresponding annotated lambda,
choosing \(z\) fresh. Map an environment to the ordered named
declarations \(\widehat\Gamma\), translating each annotation in
its older prefix. A logical context occurring as a type is translated
as a product, not flattened as an environment.

The target is exactly the Figure~3 CC of
\cref{def:p12-har-pts-specifications,thm:p12-har-cc-normalization}:
sorts \(\ast,\Box\), axiom \(\ast:\Box\), and triples
\((\ast,\ast,\ast),(\ast,\Box,\Box),
(\Box,\ast,\ast),(\Box,\Box,\Box)\).
Here \(\Box\) is the target's marker for source validity, not a
missing original source constant. Its structural metatheory and
strong normalization are independently proved, using the separately
identified 1994 primary proof
\cite[\S\S2--3]{Geuvers1994Report9450}.
\end{hbdefinition}

\begin{hblemma}{thm:p07-ch88-binding-transport}{proved=annex:p7.ch88.binding-transport}{Binding, substitution and positive beta transport}
The named encoding is injective on the scoped source grammar, modulo
consistent binder renaming, and has an indexed inverse on its image.
It preserves corrected relocation and lookup. For an object argument
\(N\) scoped over \(\Gamma\) and \(Q\) over \(\Gamma[x:A]\),
\[
 T_{\vec z}([N/x]Q)
   =T_{\vec z,z}(Q)[T_{\vec z}(N)/z].
\]
Each full compatible source beta contraction maps to one target beta
contraction at the corresponding root or annotation/body/application
position. No source annotation or type argument is erased.
This is not a bijection onto arbitrary raw CC terms, which additionally
contain \(\Box\) and ill-typed forms absent from the source grammar.
\end{hblemma}

\begin{hbtheorem}{thm:p07-ch88-full-transport}{proved=annex:p7.ch88.full-transport}{Full validity, typing and equality transport}
In the full corrected indexed presentation,
\(\Gamma\vdash\Delta\) implies
\(\widehat\Gamma\vdash_{\rm CC}T\Delta:\Box\), and
\(\Gamma\vdash M:P\) implies
\(\widehat\Gamma\vdash_{\rm CC}TM:TP\).
If \(\Gamma\vdash M\cong N\), then \(TM=_\beta TN\);
for every target pseudo type \(U\),
\(\widehat\Gamma\vdash_{\rm CC}TM:U\) holds if and only if
\(\widehat\Gamma\vdash_{\rm CC}TN:U\) holds.
Validity, typing and equality are transported jointly, including the
separately well-sorted target needed by typing conversion. The result
uses the proved structural CC metatheory, not a normalization or checking
assumption for either source variant.
\end{hbtheorem}

\begin{hblemma}{thm:p07-ch88-logical-transport}{proved=annex:p7.ch88.logical-transport}{Logical edges are positive finite beta paths}
Every derived source edge
\(\Gamma\vdash M\mathrel{\triangleright}N\) maps to a nonempty
finite compatible beta path \(TM\longrightarrow_\beta^+TN\).
The transitivity rule concatenates paths; it is not treated as a
zero-step or a single-step edge.
\end{hblemma}

\begin{hbtheorem}{thm:p07-ch88-full-normalization}{proved=annex:p7.ch88.full-normalization}{Open, fully annotated normalization for the reconstructed full calculus}
\label{thm:p07-coc-strong-normalization}
Every full source-typed object \(\Gamma\vdash M:P\), and every
valid logical-context term \(\Gamma\vdash\Delta\), is strongly
normalizing for full compatible annotated beta reduction. Arbitrary
valid open environments and contractions in annotations, types and
under binders are included. The corrected restricted variant is also
strongly normalizing for that compatible beta relation on its typable
terms, and its directed logical relation has no infinite chain starting
at a typed object or valid logical context. The proof transports into
the proved exact CC theorem; it does not assume source normalization,
closed inhabited environments, or checker completeness.
\end{hbtheorem}

\subsection{Stripped programs}

\begin{hbdefinition}{def:p07-ch88-stripping}{The source's proof-level type erasure}
For a valid environment \(\Gamma\), let \(\alpha_\Gamma\) count its
declarations of object types and let \(j_\Gamma\) inject their new
pure-program index positions into all source positions. Initially
\(\alpha_\ast=0\). Appending a logical-context annotation leaves
\(\alpha\) unchanged and adds one to every value of \(j\).
Appending an object-type annotation raises \(\alpha\) by one,
sets \(j(1)=1\), and sends each old \(j(k)\) to the new
\(j(k+1)=j(k)+1\).

For a typing \(\Gamma\vdash M:P\) with \(P\) an object,
\(\nu_\Gamma(M)\) is a pure indexed lambda term over
\(\alpha_\Gamma\) variables. At a variable \(k\), use
\(j_\Gamma^{-1}(k)\). For an application, preserve the application
when its argument has an object type, and erase that argument when
its type is a logical context. For a lambda with object-type domain,
keep an unannotated lambda over the stripped body; for a logical-context
domain, erase the abstraction and retain the stripped body. These
clauses apply only to the source's proof/functionality regime, not
arbitrarily to types or logical contexts. In the empty environment,
and only in this proof-level case, write
\(\rho(M)=\nu_\ast(M)\) for the source's closed pure-program
interpretation
\cite[\S4, pp.~105--107]{CoquandHuet1988}.
\end{hbdefinition}

\subsection{The original realization: closed programs and saturated types}
\label{sec:p07-ch88-realization}

The model of the original article's sections~4--5 is reconstructed here
in the corrected binding/equality presentations. Its proof-level
realization induction is developed in full in the accompanying annex,
with the semantic domains and argument modes made explicit. This is
distinct from the annotation-retaining interpretation used for full
annotated normalization. Programs, semantic constructor values and
realized valuations are different objects
\cite[pp.~105--110]{CoquandHuet1988}.

\begin{hbdefinition}{def:p07-ch88-native-saturation}{The source closed-program saturated sets}
Let \(\mathcal Y\) be the set of finite closed untyped lambda terms with
an inert constant \(\Omega\). Here \(\Omega\) is a symbol with no
reduction rule, not the usual divergent self-application abbreviation.
Use ordinary compatible beta reduction. A saturated set \(A\subseteq
\mathcal Y\) contains \(\Omega\), contains every
\(\Omega b_1\cdots b_n\) with strongly normalizing closed arguments,
and consists only of strongly normalizing terms. It has the source's
head-expansion closure: if \(b\) is strongly normalizing and
\((u[b/x])b_1\cdots b_n\in A\), then
\(((\lambda x.u)b)b_1\cdots b_n\in A\), whenever these whole
expressions are closed. The template \(u\) may have the one free variable
\(x\); the argument and spine terms are closed. Write \(\mathcal U\)
for all such saturated sets. No reduction-closure or general
beta-equivalence-invariance of membership is added.

For \(A\in\mathcal U\) and a total family \(F:\mathcal Y\to
\mathcal U\), define
\[
 \Pi(A,F)=\{t\in\mathcal Y:\forall b\in A,\ tb\in F(b)\}.
\]
This dependent product uses the same actual argument \(b\) in the
application and the family index. It is not an ordinary full function
space of set-theoretic functions \cite[p.~107]{CoquandHuet1988}.
\end{hbdefinition}

\begin{hblemma}{thm:p07-ch88-native-saturation}{proved=annex:p7.ch88.native-saturation}{The source closure and introduction properties}
\(\mathcal U\) is nonempty and is closed under nonempty intersections
and the displayed dependent product. If \(u[b/x]\in F(b)\) for
every \(b\in A\), with the same one-variable template \(u\), then
\(\lambda x.u\in\Pi(A,F)\). These properties do not assume source
term normalization.
\end{hblemma}

\begin{hbdefinition}{def:p07-ch88-native-functionality}{Functionality and the two ambient models}
For variant \(q\in\{r,f\}\), an object \(M\in\mathcal A_i\)
has functionality \(\varphi_q(M)=\mathcal Y\), and
\(\varphi_q(\ast)=\mathcal U\). For a logical context
\([x:A]\Delta\), the restricted model \(r\) uses all functions
\(\varphi_r(A)\to\varphi_r(\Delta)\). In the full model \(f\),
use all such functions when \(A\) is a logical context; when \(A\)
is an object, use precisely the functions
\(g:\mathcal Y\to\varphi_f(\Delta)\) satisfying
\(b=_\beta c\Rightarrow g(b)=g(c)\).
The index \(q\) explicitly qualifies the source's two choices.

For an environment \(\Gamma=[x_n:A_n]\cdots[x_1:A_1]\ast\),
the source also uses functionality for the ambient Cartesian product.
To distinguish this role from a logical-context term's function space,
write
\(\varphi_{q,\mathrm{env}}(\Gamma)=
\varphi_q(A_n)\times\cdots\times\varphi_q(A_1)\).
An ambient valuation is a tuple \(e=(a_n,\ldots,a_1)\) in that
product. A coordinate whose annotation is a logical context is a
semantic constructor value; one with an object-type annotation is a
closed program. These ambient values need not yet realize their types.
All domains are ordinary sets built by finite products/function spaces
from \(\mathcal Y\) and \(\mathcal U\). They are nonempty:
\(\Omega\) inhabits \(\mathcal Y\), the set of closed SN terms
inhabits \(\mathcal U\), and constant functions give the recursive
inhabitants, including the full beta-invariant function spaces
\cite[pp.~107--108]{CoquandHuet1988}.
\end{hbdefinition}

\begin{hbdefinition}{def:p07-ch88-native-values}{The actual two-level realization}
Write \(\rho_{q,\Gamma}(M)(e)\) for the source interpretation.
If \(M\)'s type is an object, take the stripped program
\(\nu_\Gamma(M)\) and substitute the program coordinates selected
by \(j_\Gamma\). Thus the program value is in \(\mathcal Y\) and
is independent of constructor coordinates. If \(M\)'s type is a
logical context, its constructor interpretation is defined by the following
structural clauses. A variable selects its coordinate. A constructor
lambda curries the body value over the coordinate of its domain.
A constructor application applies its function value to the actual
argument value. For a small type \([x:A]B:\ast\), put
\[
\rho_{q,\Gamma}([x:A]B)(e)=
\begin{cases}
\Pi(\rho_{q,\Gamma}(A)(e),
 b\mapsto\rho_{q,\Gamma[x:A]}(B)(e,b)),&A\text{ an object},\\
\displaystyle\bigcap_{a\in\varphi_q(A)}
 \rho_{q,\Gamma[x:A]}(B)(e,a),&A\text{ a logical context}.
\end{cases}
\]
The first line quantifies over the realized object domain; the second
over every semantic constructor value. They are different operations.

Interpretations are defined on well-sorted, role-correct endpoints.
The source-image structural typing certificates provided by
\cref{thm:p07-ch88-full-transport} certify the constructor/proof modes,
annotation scopes and function arities; the accompanying structural CC
facts use no normalization. They are not an inverse adequacy assertion.
In particular typed equality alone need not supply a typing of an
arbitrary product endpoint. Conversion uses the sorted endpoints established
by that transport. The clauses do not define an interpretation of a
logical-context term as an object inhabitant
\cite[\S5.2, pp.~108--109]{CoquandHuet1988}.
\end{hbdefinition}

\begin{hblemma}{thm:p07-ch88-native-values}{proved=annex:p7.ch88.native-values}{Defined functionalities, scoping and equality compatibility}
The constructor/proof values just defined are total on their ambient
valuation domains and belong to the functionality of their prescribed
type. Small-type values belong to \(\mathcal U\).
Relocation selects the same older coordinates and interpretation commutes
with a typed binder-removal substitution by inserting the actual value
of its argument in the valuation.

For interpretable typed-equality endpoints, functionality spaces agree.
Restricted constructor equality gives literally equal values. In the full
variant, constructor/type equality also gives literally equal values,
whereas proof equality gives beta-equal stripped programs. This last
conclusion does not assert equal saturated-set membership under arbitrary
beta equivalence. The full model's beta-invariant function requirement
is used when a constructor receives beta-equal proof arguments.
\end{hblemma}

\begin{hbdefinition}{def:p07-ch88-native-realizers}{The source realized subenvironment}
Define \(D_q(\ast)=\{()\}\). For an object-type annotation \(A\),
\[
 D_q(\Gamma[x:A])=
 \{(e,b):e\in D_q(\Gamma),\ b\in\rho_{q,\Gamma}(A)(e)\}.
\]
For a logical-context annotation \(C\), put
\(D_q(\Gamma[x:C])=D_q(\Gamma)\times\varphi_q(C)\).
Thus \(D_q(\Gamma)\subseteq\varphi_{q,\mathrm{env}}(\Gamma)\).
The dependency of an object coordinate on its older prefix is retained;
this is not the whole ambient Cartesian product. For example
\(D([A:\ast][x:A]\ast)=\{(R,b):R\in\mathcal U, b\in R\}\)
\cite[\S5.3, p.~109]{CoquandHuet1988}.
\end{hbdefinition}

\begin{hbtheorem}{thm:p07-ch88-native-soundness}{proved=annex:p7.ch88.native-soundness}{Proof-level soundness for the original realization}
For either corrected source variant with its corresponding model, if
\(\Gamma\vdash M:P\) and \(P\in\mathcal A_i^{|\Gamma|}\),
then for every \(e\in D_q(\Gamma)\),
\[
 \rho_{q,\Gamma}(M)(e)\in\rho_{q,\Gamma}(P)(e).
\]
The right side is a saturated set by the source-image sort certificate,
and the left a closed pure program. In particular this proves the original
two-premise statement when \(\Gamma\vdash P:\ast\) is also given.
This expands the source's stated structural induction; it uses neither
the full annotated normalization conclusion nor the annotation-retaining
interpretation \cite[\S5.4, p.~109]{CoquandHuet1988}.
\end{hbtheorem}

\begin{hbcorollary}{thm:p07-ch88-native-closed-image}{proved=annex:p7.ch88.native-closed-image}{The actual closed stripped-image normalization}
\label{thm:p07-ch88-source-stripped-normalization}
If \(\vdash M:P\) in either corrected source variant and
\(P\in\mathcal A_i^0\), then
\(\nu_\ast(M)=\rho_{q,\ast}(M)(())\) is strongly normalizing.
This is the source's closed program-image conclusion, not normalization
of every retained annotation of \(M\) \cite[p.~110]{CoquandHuet1988}.
\end{hbcorollary}

\begin{hbtheorem}{thm:p07-ch88-native-consistency}{proved=annex:p7.ch88.native-consistency}{The source's uninhabited proposition}
The small type \([A:\ast]A\) has no closed source proof in either
corrected variant. The inert marker belongs to the semantic interpretation
but cannot be created by beta reduction of a source-stripped program
\cite[Corollary~2, p.~110]{CoquandHuet1988}.
\end{hbtheorem}


\begin{proofboundary}{def:p07-ch88-proof-standing}{Realization, annotated normalization and extensions}
The native realization induction has its full proof in
\cref{annex:p7.ch88.native-soundness}, and its closed stripped-image
consequence is \cref{thm:p07-ch88-native-closed-image}. These develop
the original article's model and proof-level induction
\cite[\S\S4--5, pp.~105--110]{CoquandHuet1988}; they do not infer
annotation termination from erased programs. The complete annotated/open
argument is separately
\cref{thm:p07-ch88-full-transport,thm:p07-ch88-full-normalization},
using the independently supplied exact CC proof. The article's invocation
of full constructed-term normalization on pp.~105 and 112 and its
reference~15 retain their distinct attribution. The complete checking
result in \cref{thm:p07-ch88-check-normal-complete} concerns the separate
no-conversion normal-type variant. Ordinary-full constructor root
preservation has its actual source casts in
\cref{thm:p07-ch88-ordinary-root}; proof-valued conversions in the
certified head-excursion family are handled by
\cref{thm:p07-ch88-object-normalize,thm:p07-ch88-object-injection}.
These source constructions do not establish arbitrary headedness or
global ordinary-full checking. The remaining exhaustive typed-equality
dependency is stated in \cref{def:p07-ch88-object-boundary}.

Section~7's constants and local let are abbreviation facilities; no
extra delta generator has been added here. Section~8's primitive
recursors, universe proposals and Type:Type are different extensions,
not consequences of the base normalization theorem. The article's
warning that some proposed universe rules lose normalization is retained
as a limit, not a claim about every stratified universe system
\cite[pp.~105,112--117]{CoquandHuet1988}.
\end{proofboundary}

\section{Original logical rank and the normal-type checking variant}
\label{sec:p07-ch88-checking}

This section uses the corrected indexed binding and the full/restricted
source rules of \cref{sec:p07-ch88}. The original checking discussion is
on pp.~104--105; the separate no-conversion, normal-type presentation is
\S6.1, pp.~111--112 \cite{CoquandHuet1988}. The latter has a complete
checker below, with a direct adequacy proof for its declared domain.
It is not silently identified with the full calculus on unchanged raw
terms. All additional written rules are on
Sheet~\ref{sheet:p07-ch88-normal-types}.

Named expressions in the examples denote the indexed inverse of the
role-sensitive encoding \(T\), with a fixed external variable vector.
This is the proved raw binding correspondence of
\cref{thm:p07-ch88-binding-transport}, not an assumed inverse of target
typing. In particular the named example variables are actual bound or
external references, not the dummy letters of indexed binder labels.

\begin{hbmethod}{def:p07-ch88-source-checking}{Logical normalization and the distinct checking presentations}
For a constructed logical redex \(((\lambda x:P)M)N\), assign
domain complexity \(\delta(P)\), and collect these complexities
in the domain-rank multiset of a construction. Rank invariance is
\cref{thm:p07-ch88-check-rank}; this multiset nevertheless need not
decrease under an innermost logical contraction, as the typed example
in \cref{thm:p07-ch88-check-domain-rank} proves. The logical relation
instead terminates on constructed terms by
\cref{thm:p07-ch88-logical-transport,thm:p07-ch88-full-normalization}.

Realization proves semantic membership and closed erased-program
termination, not a type-inference algorithm. Source generation and
substitution are supplied below. The explicit algorithm and its
soundness/completeness proof concern the p.~112 normal-type variant;
its normalized declarations/results and absence of conversion are
part of its specification. Ordinary-full checking has the separate
typed-equality dependency at the end of this section
\cite[pp.~104--105,112]{CoquandHuet1988}.
\end{hbmethod}

\begin{hblemma}{thm:p07-ch88-check-generation}{proved=annex:p7.ch88.check-generation}{Source context regularity and conversion-stripping generation}
Every derived validity, typing or typed-equality judgment in either
corrected source variant has a valid environment. Its declarations have
their actual older-prefix certificates: a logical-context annotation is
valid there, or an object annotation has type \(\ast\) there.

If \(\Gamma\vdash M:P\), stripping final typing conversions exposes
the common rule for the literal subject constructor. A variable's initial
type is its lookup; an object product has initial type \(\ast\) and a
body premise typed \(\ast\); a lambda has a body typing and the
corresponding literal product type; an application has function/argument
premises and initial type \([N/x]B\). The initial type is joined to
\(P\) by the actual finite chain of source typing-conversion premises,
allowing a zero-length chain. This statement does not assume that every
raw equality endpoint is itself a sorted type.
\end{hblemma}

\begin{hbtheorem}{thm:p07-ch88-check-substitution}{proved=annex:p7.ch88.check-substitution}{Actual source thinning and dependent substitution}
In either corrected ordinary source variant, a judgment may be thinned
to a valid environment containing its declarations in the same relative
order, with the same annotations and fresh additional declarations.
The indexed conclusion uses the corresponding reference relocation.
If \(\Gamma[x:A];\Theta\vdash J\) and \(\Gamma\vdash N:A\),
where \(N\) is a source object, then
\(\Gamma;([N/x]\Theta)\vdash[N/x]J\). This covers validity, typing
and typed equality; every suffix annotation is substituted in its own
older prefix. No beta subject-reduction or inverse CC typing premise
is needed.
\end{hbtheorem}

\begin{hblemma}{thm:p07-ch88-check-rank}{proved=annex:p7.ch88.check-rank}{The source rank-invariance lemma}
For the source rank \(\delta\) of
\cref{def:p07-ch88-logical-reduction}, an object-argument substitution
preserves the rank of every logical-context term. Source typed equality
preserves the object/logical-context class and \(\delta\), in either
corrected variant. Rank invariance alone does not imply a decreasing
redex-domain multiset; the next result separates those properties.
\end{hblemma}

\begin{hbtheorem}{thm:p07-ch88-check-domain-rank}{proved=annex:p7.ch88.check-domain-rank}{A typed boundary for the domain-rank measure}
In the valid open environment \(A:\ast,a:A,C:\ast\), put
\[
 L=(\lambda x:A)(\lambda B:\ast)B,\qquad
 M=(La)C,\qquad M'=((\lambda B:\ast)B)C.
\]
Both \(M\) and \(M'\) have source type \(\ast\), even in the
corrected restricted variant. An innermost permitted logical contraction
gives \(\Gamma\vdash M\triangleright M'\).
The domain-rank multiset changes from \(\{0\}\) to
\(\{1\}\). Consequently that multiset is not a
decreasing measure for these rules. The already proved normalization
theorem is independent of that measure.
\end{hbtheorem}

\begin{hblemma}{thm:p07-ch88-check-raw-reflection}{proved=annex:p7.ch88.check-raw-reflection}{Raw beta reflection and unique normal forms}
A target compatible beta step starting at an encoded scoped source term
reflects to exactly one scoped source beta step. Every target reduction
path from that image remains in the raw source image. Therefore the full
compatible raw source relation is confluent. On any source image supplied
with a CC typing certificate, full normalization terminates and has a
unique indexed normal form. No inverse of source typing or typed equality
is part of this assertion.
\end{hblemma}

\begin{hbdefinition}{def:p07-ch88-check-normal-variant}{The exact no-conversion normal-type variant}
Write \(\Gamma\vdash_n\Delta\) and \(\Gamma\vdash_n M:P\) for
the variant on Sheet~\ref{sheet:p07-ch88-normal-types}. It has the
eight p.~112 construction rules, normalized declaration annotations and
normalized application result types; it has no typing-conversion or
typed-equality rule. The notation \(\mathcal N(E)\) is the result of
a fixed full compatible beta-normalizing strategy, when that strategy
terminates, not a total function on arbitrary raw terms. The theorem
below proves that its calls at derived rule premises terminate.

Call a source expression \emph{annotation-normal} if every product or
lambda domain annotation is in full beta normal form, and its other
children are annotation-normal recursively. This is a finite syntactic
test; it does not normalize a possibly divergent input. Every derivable
environment annotation in this variant is fully normal.
\end{hbdefinition}

\calculussheet{sheet:p07-ch88-normal-types}{Coquand--Huet 1988 normal-type variant}{Corrected indexed binding and the eight p.~112 construction rules. Full beta normal forms are used only at the declared normalized annotations/results; no conversion or typed-equality rule. This is a separate presentation, not raw-beta subject-preserving.}
\hbrule{rule:p07-ch88-n-empty}{Empty validity}{}{\vdash_n\ast}{\none}
\hbrule{rule:p07-ch88-n-kind-extension}{Normalized logical declaration}{\Gamma\vdash_n\Delta}{\Gamma[x:\mathcal N(\Delta)]\vdash_n\ast}{The normalizer is defined on this typed kind by the direct CC transport proof.}
\hbrule{rule:p07-ch88-n-type-extension}{Normalized object-type declaration}{\Gamma\vdash_n M:\ast}{\Gamma[x:\mathcal N(M)]\vdash_n\ast}{The normalizer retains annotations and includes every compatible beta position.}
\hbrule{rule:p07-ch88-n-kind-product}{Logical product}{\Gamma[x:A]\vdash_n\Delta}{\Gamma\vdash_n[x:A]\Delta}{The premise's valid environment forces its declaration annotation to be normal.}
\hbrule{rule:p07-ch88-n-type-product}{Object product}{\Gamma[x:A]\vdash_n B:\ast}{\Gamma\vdash_n[x:A]B:\ast}{The object body is in the extended indexed scope.}
\hbrule{rule:p07-ch88-n-variable}{Projection}{\Gamma\vdash_n\ast}{\Gamma\vdash_n l:\Gamma/l}{$1\le l\le|\Gamma|$; the actual corrected lookup preserves normality.}
\hbrule{rule:p07-ch88-n-abstraction}{Abstraction}{\Gamma[x:A]\vdash_n M:P}{\Gamma\vdash_n(\lambda x:A)M:[x:A]P}{The domain is the literal annotation appearing in the premise environment.}
\hbrule{rule:p07-ch88-n-application}{Application with normalized result}{\Gamma\vdash_n M:[x:A]B\qquad\Gamma\vdash_n N:A}{\Gamma\vdash_n(MN):\mathcal N([N/x]B)}{The object argument and binder-removal substitution retain their source class; target sorting makes the result normalizer terminate.}


\begin{hbtheorem}{thm:p07-ch88-check-normal-transport}{proved=annex:p7.ch88.check-normal-transport}{All-rule transport of the normal-type variant}
Every normal-variant validity maps to \(T\Delta:\Box\), and every
normal-variant typing maps to \(TM:TP\), in exact two-sort CC and its
valid encoded context. All rule calls to \(\mathcal N\) are defined.
The subject is annotation-normal, the environment has normal annotations,
and every resulting source type is fully normal.
The target is exactly \(\ast:\Box\) with its four declared product
triples, not a universe extension or a proof-erased calculus.
\end{hbtheorem}

\begin{hbdefinition}{def:p07-ch88-check-algorithm}{A total algorithm for the normal-type variant}
Reject malformed or out-of-scope input before doing any normalization.
Validate a declaration list oldest-first, retaining the derivation and
CC certificate of each accepted prefix. In a checked prefix use these
mutually structural operations:
\begin{enumerate}
\item \(\mathrm{Sort}(A)\) first requires that \(A\) is syntactically
      fully normal. For a logical context it calls \(\mathrm{Valid}(A)\).
      For an object it calls \(\mathrm{Infer}(A)\) and requires the
      returned normal type to be literally \(\ast\).
\item \(\mathrm{Valid}(\ast)\) succeeds. For \([x:A]\Delta\),
      check \(\mathrm{Sort}(A)\), append its checked declaration,
      and check \(\mathrm{Valid}(\Delta)\) in that prefix.
\item \(\mathrm{Infer}(k)\) returns the checked context lookup.
      For an object product \([x:A]B\), check \(\mathrm{Sort}(A)\),
      infer \(B\) in the checked extension, and require its type
      literally \(\ast\); return \(\ast\).
\item For \((\lambda x:A)M\), check \(\mathrm{Sort}(A)\), infer
      the body type \(P\) in the checked extension, and return
      \([x:A]P\).
\item For \((MN)\), infer the normal function type. It must be a
      literal product \([x:A]B\). Infer the normal argument type
      and require literal equality with \(A\). Return
      \(\mathcal N([N/x]B)\). This last call is made only after the
      preceding successes have produced its target sort certificate.
\end{enumerate}
Any failed structural or equality check returns rejection. Each
recursive syntax call visits a proper child of its current original
expression, without rechecking the whole environment. Context validation
is a finite prefix loop. Lookup and comparison use the corrected indexed
operations, so equality of outputs is literal indexed equality, not
agreement of arbitrarily chosen binder names.
\end{hbdefinition}

\begin{hbtheorem}{thm:p07-ch88-check-normal-complete}{proved=annex:p7.ch88.check-normal-complete}{Soundness, completeness and termination of that checker}
The algorithm is total on finite scoped inputs and decides exactly the
normal variant's validity and typing judgments. On success it constructs
an actual derivation and, for an object, its unique literal normal type.
Thus supplied-type checking is equality with that returned type.

More precisely, in a checked annotation-normal environment, if
\(\widehat\Gamma\vdash_{\rm CC}TM:U\) for an annotation-normal
source object \(M\), the algorithm succeeds with a type \(P\) and
\(TP=_\beta U\). A source logical context typed in CC at \(\Box\)
is accepted by \(\mathrm{Valid}\). This is proved by structural
generation in CC, including the application product-head and normalized
domain cases; it is not an assumed inverse to the full source bridge.
Any encoded source environment which is CC-valid and has normal
annotations is likewise accepted by context validation.

Every derived \(\Gamma\vdash_n M:P\) has the actual source type
certificate \(\Gamma\vdash_n P:\ast\), if \(P\) is an object,
or \(\Gamma\vdash_n P\), if it is a logical context. It also has
\(\Gamma\vdash_n\mathcal N(M):P\). Preservation here is to the
whole normal form, not to every intermediate raw beta contractum.
\end{hbtheorem}

\begin{hbtheorem}{thm:p07-ch88-check-normal-step}{proved=annex:p7.ch88.check-normal-step}{Why normal types are a separate presentation}
In the valid normal environment \(A:\ast,a:A\), let
\[
 K=(\lambda x:A)A,\qquad
 V=((\lambda F:[x:A]\ast)(\lambda y:Fa)y)K.
\]
The normal variant derives \(V:[y:A]A\). Its raw root contractum
\((\lambda y:Ka)y\) has a nonnormal annotation and is not typable
in that variant. Its whole normal form \((\lambda y:A)y\) is typable.
Thus this variant is not raw-beta subject-preserving, and replacing a
raw full-calculus judgment by an unchanged normal-variant judgment is
not an adequacy theorem.
\end{hbtheorem}

\begin{proofboundary}{def:p07-ch88-check-full-boundary}{The remaining ordinary-full source checking dependency}
\label{def:p07-ch88-source-checking-boundary}
For an annotation-normal full-source environment and subject, the proved
forward full transport followed by the new checker gives a normal-variant
derivation at the normal form of its old type. The converse ordinary-full
judgment and the checker on arbitrary full-source annotations still need
ordinary-full source type correctness, type preservation and completeness of typed source
equality on beta-convertible sorted endpoints. Raw image reflection alone
does not supply those derivations. They are not assumed here, and no
complete ordinary-full or restricted typed-equality checker is claimed.
The proved complete algorithm belongs to the explicitly declared p.~112
no-conversion variant.
\end{proofboundary}

\section{Ordinary-full source conversions and constructor redexes}
\label{sec:p07-ch88-ordinary-checking}

These are further source-rule constructions, not an identification of the
ordinary full calculus with its normal-type variant. Only the corrected
ordinary full rules on Sheet~\ref{sheet:p07-ch88-full} are used. In
particular an object declaration requires an actual derivation of its
annotation at \(\ast\); target typing alone does not license it.
Named expressions below use the already proved raw indexed correspondence.
The original ordinary rules are \cite[pp.~99--102]{CoquandHuet1988}.

\begin{hblemma}{thm:p07-ch88-ordinary-star}{proved=annex:p7.ch88.ordinary-star}{Literal-star rigidity and its generation consequences}
For any scoped source expression \(E\), if \(TE=_\beta\ast\), then
\(E\) is literally \(\ast\). Thus source equality cannot convert
\(\ast\) to a different expression. No ordinary source lambda has
type \(\ast\). A typing \(\Gamma\vdash(MN):\ast\) has actual
generation premises \(\Gamma\vdash M:[x:A]\ast\) and
\(\Gamma\vdash N:A\). This last assertion does not supply an independent
source sort certificate for arbitrary inferred \(A\).
\end{hblemma}

\begin{hbtheorem}{thm:p07-ch88-ordinary-context}{proved=annex:p7.ch88.ordinary-context}{Context conversion with actual source annotation certificates}
Suppose \(A,A'\) each have their actual annotation certificate in
\(\Gamma\): object annotations are typed \(\ast\), and logical
annotations are valid logical contexts. Suppose also
\(\Gamma\vdash A\cong A'\). For every ordinary-full validity,
typing or equality judgment in \(\Gamma[x:A];\Theta\), one can
construct the same named judgment in \(\Gamma[x:A'];\Theta\).
The target suffix is valid; its dependent annotations are not discarded
or normalized. Indexed decoding gives the actual corresponding judgments.
No beta-equivalence-to-source-equality reflection is a premise.
\end{hbtheorem}

\begin{hbtheorem}{thm:p07-ch88-ordinary-equal-substitution}{proved=annex:p7.ch88.ordinary-equal-substitution}{Equality under two source-equal substitutions}
Suppose \(\Gamma\vdash N:A\), \(\Gamma\vdash N':A\), and
\(\Gamma\vdash N\cong N'\). If
\(\Gamma[x:A];\Theta\vdash E:P\), or
\(\Gamma[x:A];\Theta\vdash E\) is a validity judgment, then
\[
 \Gamma;([N/x]\Theta)\vdash [N/x]E\cong[N'/x]E.
\]
The right expression is compared in the first substituted environment,
as required by the full product/lambda congruence rules. This theorem
uses full typed reflexivity and simultaneous application congruence;
it is not advertised for arbitrary proof arguments in the restricted
variant.
\end{hbtheorem}

\begin{hbdefinition}{def:p07-ch88-ordinary-product-headed}{Product-headed equality proofs}
A product-headed source equality proof has products at both endpoints
of every outer symmetry/transitivity node. Its outer leaves are product
congruence or reflexivity of a typed object product or valid logical
product. The annotation/body equality subproofs may be arbitrary full
source proofs. Outer lambda, application and beta generators are excluded.
This is a property of an actual finite proof, not an assumption that
every object-product equality can already be transformed into that form.
\end{hbdefinition}

\begin{hbtheorem}{thm:p07-ch88-ordinary-product-injection}{proved=annex:p7.ch88.ordinary-product-injection}{Actual source product injection for the headed regime}
Suppose \(A\) has its actual source annotation certificate in
\(\Gamma\), and a product-headed proof gives
\(\Gamma\vdash[x:A]P\cong[x:B]Q\). Then one can construct
\[
 \Gamma\vdash A\cong B,\qquad
 \Gamma[x:A]\vdash P\cong Q.
\]
No certificate for the final \(B\) is assumed. If both endpoint
products are logical-context terms, every source equality proof between
them is automatically product-headed. Hence the conclusion is
unconditional for that kind-product case. It is not a claimed general
injection theorem for object products with arbitrary beta excursions.
\end{hbtheorem}

\begin{hbtheorem}{thm:p07-ch88-ordinary-root}{proved=annex:p7.ch88.ordinary-root}{Root subject reduction at every ordinary constructor type}
Let \(R=((\lambda x:A)M)N\) and \(C=[N/x]M\).
If \(\Gamma\vdash R:\Delta\), where \(\Delta\) is a logical
context, then actual ordinary-full derivations give
\[
 \Gamma\vdash C:\Delta,\qquad\Gamma\vdash R\cong C.
\]
The generation argument constructs the actual source domain
equality and casts the argument to the literal lambda annotation.
For an arbitrary result type, the same conclusion holds when the actual
function-type conversion chain exposed by generation has product-headed
equality proofs. With literal argument typing \(\Gamma\vdash N:A\)
and body typing \(\Gamma[x:A]\vdash M:P\), both endpoints always
have the actual principal type \([N/x]P\); any explicitly supplied
outer conversion chain can be reapplied to the contractum.
\end{hbtheorem}

\begin{hbtheorem}{thm:p07-ch88-ordinary-application}{proved=annex:p7.ch88.ordinary-application}{Dependent application types under an actual equal argument}
Suppose \(\Gamma[x:A]\vdash B:\ast\), or that logical \(B\)
is valid there; suppose \(\Gamma\vdash F:[x:A]B\) and the same
three actual argument judgments as in
\cref{thm:p07-ch88-ordinary-equal-substitution}. Then
\[
 \Gamma\vdash[N/x]B\cong[N'/x]B,\qquad
 \Gamma\vdash FN':[N/x]B.
\]
The first is a constructed source equality, not merely a target beta
comparison. This supplies the codomain-transport step once the two
arguments' source typings at the same domain have been obtained.
\end{hbtheorem}

\begin{hbtheorem}{thm:p07-ch88-ordinary-example}{proved=annex:p7.ch88.ordinary-example}{Nonnormal declarations and an annotation contraction with the same old type}
In the ordinary valid environment
\[
 \Gamma=A:\ast,\ F:[u:A]\ast,\ a:A,\ f:[u:A]Fu,
\]
put \(n=((\lambda v:A)v)a\), \(B=Fn\), \(B'=Fa\), and
\(T=[h:B]B\). Actual full-source derivations give
\[
 n:A,\quad n\cong a,\quad B:\ast,\quad B':\ast,
 \quad B\cong B',\quad fn:B,\quad fa:B.
\]
The nonnormal declaration \(\Gamma[h:B]\) is valid. Changing it
to \(\Gamma[h:B']\) preserves the old judgment \(h:B\).
Furthermore \(t=(\lambda h:B)h\) contracts through its annotation
to \(t'=(\lambda h:B')h\), and both have the actual old type \(T\)
with \(\Gamma\vdash t\cong t'\).
The no-conversion normal-type variant cannot type the old nonnormal
annotation or retain the nonnormal old type \(T\).
\end{hbtheorem}

\begin{proofboundary}{def:p07-ch88-ordinary-boundary}{The remaining ordinary proof-valued conversion case}
General ordinary-full type correctness, sorted-beta equality reflection
and compatible subject reduction are not asserted complete here.
Constructor-valued root reduction has its source product injection and
actual cast. For a proof-valued redex, the function type is an object
product, whose equality proof can leave the product head through typed
beta/application generators. The next section's certified transformation
handles a specified recursive family of these genuinely unheaded proofs,
and extends root preservation to their conversion chains
(\cref{thm:p07-ch88-object-normalize,thm:p07-ch88-object-injection}).
It does not establish exhaustive head elimination. Neither one-way CC
typing, normal-type checking nor successful nonnormal examples supplies
global sorted reflection or an arbitrary-annotation checker. The final
scope boundary is \cref{def:p07-ch88-object-boundary}.
\end{proofboundary}

\section{Certified elimination of object-head excursions}
\label{sec:p07-ch88-object-product}

This is a proof transformation for a specified family of actual ordinary-
full equality proofs. It uses the original guards on
Sheet~\ref{sheet:p07-ch88-full}, the proved source context conversion and
substitution constructions, and no inverse typing theorem. Its certificate
condition is explicit syntax of a finite supplied proof, not the assertion
that the desired headed proof already exists. It does not yet cover every
object-product equality proof \cite[pp.~99--102]{CoquandHuet1988}.

\begin{hbdefinition}{def:p07-ch88-object-height}{Two different finite heights}
For a strongly normalizing scoped source term \(E\), let
\(\eta_\beta(E)\) be the maximum length of a full compatible beta
reduction from it, including annotation and body positions. Finite branching
and strong normalization make this a natural number. For an original
finite source proof \(d\), let \(\operatorname{ht}(d)\) be its tree
height, counting all typing, validity and equality premises. Generated
proofs may have greater height; the transformation below never recurses
on those generated proofs.
\end{hbdefinition}

\begin{hblemma}{thm:p07-ch88-object-landing}{proved=annex:p7.ch88.object-landing}{Beta-product landings have actual source small-type guards}
Suppose the actual full beta premises are
\(\Gamma[x:A]\vdash M:P\) and \(\Gamma\vdash N:A\), and
\(C=[N/x]M\) is literally an object product. Then \(P\) is literally
\(\ast\), and source derivations give \(M:\ast\) in its extended
context, \(C:\ast\), and \(((\lambda x:A)M)N:\ast\).

More generally, take two primitive beta landings with their actual
premises and the same object head \(H\), where \(H\) is product
or lambda. If
\(\Gamma[x:A]\vdash M\cong M'\), their templates are either
both literally \(H\)-headed or both the selected variable \(x\).
Mixed cases are impossible. In this shape assertion the right landing
may use annotation \(A'\), with its actual original beta premises.
\end{hblemma}

\begin{hbtheorem}{thm:p07-ch88-object-asymmetric}{proved=annex:p7.ch88.object-asymmetric}{Asymmetric equal substitution needs no right argument typing}
If \(\Gamma\vdash N:A\), \(\Gamma\vdash N\cong N'\), and
\(N'\) is a scoped object, then for a supplied typing or validity
\(\Gamma[x:A];\Theta\vdash E:P\) or \(\Gamma[x:A];\Theta\vdash E\),
\[
 \Gamma;([N/x]\Theta)\vdash[N/x]E\cong[N'/x]E.
\]
No derivation of \(N':A\) is needed. If \(E\) is itself an object
product or lambda with a supplied typing, this construction has that
same head at its outer congruence leaf; final typing conversions in the
template proof do not change its subject.
\end{hbtheorem}

\begin{hbdefinition}{def:p07-ch88-object-diamond}{An original guarded beta--congruence--beta diamond}
Put
\[
 R_0=((\lambda x:A)M)N,\quad C_0=[N/x]M,
 \qquad R_1=((\lambda x:A')M')N',\quad C_1=[N'/x]M'.
\]
The supplied equality proof follows \(C_0\cong R_0\cong R_1\cong C_1\).
Its first and last edges are symmetry of primitive beta and primitive
beta, with their actual body/argument typing premises. Its middle edge
is full application congruence whose function premise is one primitive
full lambda congruence, with actual premises
\(A\cong A'\), \(M\cong M'\) in \(\Gamma[x:A]\), and the
stated left body typing. Its argument premise is \(N\cong N'\).
Both contracta are literal objects with the same head \(H\).
These are supplied source proof nodes and original guards, not inferred
source certificates from target typing. Fresh binder names are fixed
before indexed decoding.
\end{hbdefinition}

\begin{hbtheorem}{thm:p07-ch88-object-commute}{proved=annex:p7.ch88.object-commute}{A genuine cut commutation with the original guards}
Every such diamond has an actual source replacement
\[
 \Gamma\vdash C_0\cong[N/x]M'\cong C_1.
\]
The first edge is typed single substitution of the actual body equality;
the second is equal substitution in the right body, whose required first-
context typing is constructed from the landing guard by source context
conversion. For product landings both contracta are actually typed
\(\ast\). The two beta contractions strictly decrease
\(\eta_\beta(R_0)+\eta_\beta(R_1)\) when their endpoints alone
are compared with \(C_0,C_1\). This is not claimed to bound the height
of the new intermediate \([N/x]M'\).
\end{hbtheorem}

\begin{hbdefinition}{def:p07-ch88-object-supported}{Syntactic excursion certificates and headed proofs}
Fix \(H\) to be product or lambda. An \(H\)-headed proof has that
literal head at every outer symmetry/transitivity endpoint; its outer
leaves are that constructor's full congruence or typed reflexivity.
Annotation/body subproofs are opaque there.

An excursion certificate on an original equality proof with \(H\)-headed
object endpoints is built from the following syntactic cases:
\begin{enumerate}
\item an already \(H\)-headed proof;
\item symmetry or transitivity of certified original subproofs, with an
      \(H\)-headed intermediate in the transitivity case;
\item an exact cancelling beta round trip with identical redex;
\item the guarded diamond just defined. If its templates are literally
      \(H\)-headed, its original body-equality subproof must carry a
      certificate for \(H\). If both templates are the selected variable,
      its original argument-equality subproof must carry that certificate.
\end{enumerate}
Each recursive child is a proper subproof of the ORIGINAL proof. No
certificate is stipulated for a newly generated equality. Recognition of
these cases and the raw substitution/head tests is finite on the supplied
proof; this definition makes no exhaustion claim for all valid proofs.
\end{hbdefinition}

\begin{hbtheorem}{thm:p07-ch88-object-normalize}{proved=annex:p7.ch88.object-normalize}{A terminating headed-proof transformation for the certified family}
Every certified proof can be transformed to an actual \(H\)-headed
source proof with the same context and endpoints. The algorithm recurses
strictly on \(\operatorname{ht}(d)\) of its original equality proof,
including the diamond's original body or argument subproof. It never
uses the possibly larger height of a generated substitution/context-
conversion proof as a recursive input.
\end{hbtheorem}

\begin{hbtheorem}{thm:p07-ch88-object-injection}{proved=annex:p7.ch88.object-injection}{Product injection and root preservation for genuine unheaded certificates}
For a certified object-product equality
\(\Gamma\vdash[x:A]P\cong[x:B]Q\), with the actual annotation
certificate for \(A\), construct
\(\Gamma\vdash A\cong B\) and
\(\Gamma[x:A]\vdash P\cong Q\).
The ordinary root subject-reduction construction consequently applies to
proof-valued redexes whose actual function-type conversion chains have
these certificates, not just already-headed chains.

For an \(H\)-headed lambda equality with an actual first-annotation
certificate, the analogous conclusions are the domain equality and the
body equality in that first annotation's context. No right-domain source
certificate is silently assumed.
\end{hbtheorem}

\begin{hbtheorem}{thm:p07-ch88-object-height-obstruction}{proved=annex:p7.ch88.object-height-obstruction}{A typed obstruction to the naive outer-path height measure}
In \(\Gamma=A:\ast,C:\ast,a:A\), let
\[
 I_A=(\lambda v:A)v,\quad n_0=a,\quad n_{j+1}=I_A n_j,
 \quad D(x)=((\lambda p:A)C)x,
\]
and put
\[
 M_0=[y:C]C,\qquad M_1(x)=[y:D(x)]D(x),
 \quad R_0=((\lambda x:A)M_0)n_3,
 \quad R_1=((\lambda x:A)M_1(x))a.
\]
There is a genuine guarded product diamond, with product-headed body
proof and argument equality \(n_3\cong a\). Its original outer path
has beta heights \(0,4,3,2\). The commuted intermediate \(M_1(n_3)\)
has height \(8\), so the new outer path has heights \(0,8,2\).
Both outer-path maximum and sum increase. This does not disprove head
elimination: the certified algorithm terminates by original proof height
and returns headed proofs in this example.
\end{hbtheorem}

\begin{proofboundary}{def:p07-ch88-object-boundary}{The remaining exhaustion and general-guard problem}
The certificate family includes real nested unheaded beta diamonds, but
does not yet cover an arbitrary equality between object products. A
middle function equality need not be one primitive lambda congruence;
several application/conversion steps may expose its lambda head. No
decreasing measure for eliminating all such shapes has been proved here.
The explicit height obstruction excludes one naive proposed measure, not
the theorem's possibility. Global ordinary-full type correctness, general
sorted reflection, compatible subject reduction and arbitrary-annotation
checking remain separate unsupplied conclusions.
\end{proofboundary}

